% Irodov 2.3, source PDF pp.80--81, printed pp.71--72.
% Problem 74 source answer appears to have a factor-of-ten numeric error.
\thermokineticproblem{73}{71}
一混合气由 $\nu_1$ mol 单原子理想气体和
$\nu_2$ mol 刚性双原子理想气体组成。
求混合物的热容比 $\gamma$。
\exsolution
单原子分子的定容、定压摩尔热容为
$3R/2$、$5R/2$；刚性双原子分子对应
$5R/2$、$7R/2$。
混合气总热容为
\[
C_{V,\Sigma}=\frac R2(3\nu_1+5\nu_2),\qquad
C_{p,\Sigma}=\frac R2(5\nu_1+7\nu_2).
\]
因此
\[
\boxed{\gamma=
\frac{C_{p,\Sigma}}{C_{V,\Sigma}}
=\frac{5\nu_1+7\nu_2}{3\nu_1+5\nu_2}}.
\]
\excheck 仅剩单原子气体时 $\gamma=5/3$，
仅剩刚性双原子气体时 $\gamma=7/5$，
与独立组分一致。

\thermokineticproblem{74}{71}
一个装有氮气的绝热密闭容器
最初温度为 $t=27^\circ\mathrm C$，
整体以速率 $v=100\,\mathrm{m/s}$ 平动。
若容器突然停止，且气体的整体平动动能
转化为内能，求压强变化的方向及百分比。
\exsolution
设氮气质量 $m$、摩尔质量
$M=28\times10^{-3}\,\mathrm{kg/mol}$。
在刚性双原子气体近似中
$C_V=5R/2$，所以
\[
\frac12mv^2=\frac mM C_V\Delta T
\quad\Longrightarrow\quad
\Delta T=\frac{Mv^2}{5R}.
\]
因为容器体积及气体摩尔数不变，
$p\propto T$，故
\[
\boxed{\frac{\Delta p}{p}
=\frac{\Delta T}{T}
=\frac{Mv^2}{5RT}>0}.
\]
代入 $T\simeq300\,\mathrm K$：
\[
\boxed{100\frac{\Delta p}{p}
=\frac{100(0.028)(100)^2}{5(8.314)(300)}
\simeq2.24\%}.
\]
\excheck 复核量纲与算术：
$\Delta T=Mv^2/(5R)\simeq6.74\,\mathrm K$，
故 $100\Delta T/T\simeq2.24\%$，
与本版书末给出的约 $2.2\%$ 一致。
此前译稿误将百分比算小10倍，
现已更正；\emph{原书这一项无数值勘误}。

\thermokineticproblem{75}{71}
在温度 $t=17^\circ\mathrm C$ 下，求
(a) 氧气分子的均方根速率
与单分子的平均平动动能；
(b) 空气中悬浮的水滴
（直径 $d=0.10\,\mathrm{\mu m}$）
的布朗平动均方根速率。
\exsolution
(a) 温度为 $T=290.15\,\mathrm K$，
氧气摩尔质量
$M_{\mathrm O_2}\simeq0.0320\,\mathrm{kg/mol}$。
由能量均分：
\[
\boxed{v_{\rm rms}=\sqrt{\frac{3RT}{M_{\mathrm O_2}}}
\simeq0.47\,\mathrm{km/s}},\qquad
\boxed{\bar\varepsilon_{\rm trans}=\frac32k_BT
\simeq6.0\times10^{-21}\,\mathrm J}.
\]
(b) 将水滴视为密度
$\rho_w\simeq10^3\,\mathrm{kg/m^3}$
的刚性球体，
其质量
$m_d=(\pi/6)\rho_wd^3$。
悬浮平动的三个速度自由度平均动能为
$3k_BT/2$，因此
\[
\boxed{v_{d,\rm rms}=
\sqrt{\frac{3k_BT}{m_d}}
=\sqrt{\frac{18k_BT}{\pi\rho_wd^3}}
\simeq0.15\,\mathrm{m/s}}.
\]
\excheck 水滴的平动热运动能量和单个氧分子的平均平动能
同为 $3k_BT/2$，
但两者质量相差巨大，因此速率不同。

\thermokineticproblem{76}{71}
由刚性双原子分子构成的理想气体
发生可逆绝热膨胀。
若分子均方根速率下降
$\eta=1.50$ 倍，求气体体积增大多少倍。
\exsolution
分子均方根速率
$v_{\rm rms}\propto\sqrt T$，
所以末温
$T_2/T_1=\eta^{-2}$。
刚性双原子气体
$i=5$、$\gamma=1+2/i=7/5$。
可逆绝热关系
$TV^{\gamma-1}=\text{常数}$，故
\[
\left(\frac{V_2}{V_1}\right)^{\gamma-1}
=\frac{T_1}{T_2}=\eta^2.
\]
因此
\[
\boxed{\frac{V_2}{V_1}
=\eta^{2/(\gamma-1)}
=\eta^5
=(1.50)^5\simeq7.6}.
\]
\excheck 温度和均方根速率在绝热膨胀中都降低，
体积增大，符合物理方向。

\thermokineticproblem{77}{71}
容器中有 $m=15\,\mathrm g$ 的氮气，
初温 $T_1=300\,\mathrm K$。
为使分子均方根速率增大到原来的
$\eta=2.0$ 倍，在\emph{等容}条件下
需向气体传递多少热？
\exsolution
分子均方根速率与温度平方根成正比，
因此
$T_2/T_1=\eta^2=4$、
$\Delta T=(\eta^2-1)T_1=900\,\mathrm K$。
刚性双原子氮气的定容摩尔热容
$C_V=5R/2$，
摩尔数
$\nu=m/M=15/28\,\mathrm{mol}$。
等容无气体外功，第一定律给
\[
Q=\nu C_V\Delta T
=\frac mM\frac{5R}{2}T_1(\eta^2-1).
\]
故
\[
\boxed{Q\simeq
\frac{15}{28}\frac{5(8.314)}{2}(300)(3)
\simeq1.0\times10^4\,\mathrm J=10\,\mathrm{kJ}}.
\]
\excheck 题面未写升温过程的具体热容模型，
本稿采用\emph{刚性双原子近似}。
若高温显著激发振动，自由度与热容会改变。

\thermokineticproblem{78}{72}
由刚性双原子分子组成的气体
温度为 $T=300\,\mathrm K$。
单个分子的转动惯量为
$I=2.1\times10^{-39}\,\mathrm{g\,cm^2}$。
求分子旋转的角均方根速率。
\exsolution
刚性线性双原子分子有两项独立转动动能，
因此平均转动动能为 $k_BT$。
若转动角速度模为 $\omega$，
旋转动能 $\varepsilon_{\rm rot}=I\omega^2/2$，
故
\[
\frac12I\overline{\omega^2}=k_BT,\qquad
\boxed{\omega_{\rm rms}=\sqrt{\frac{2k_BT}{I}}}.
\]
转换惯量单位：
$1\,\mathrm{g\,cm^2}=10^{-7}\,\mathrm{kg\,m^2}$，
因此
$I=2.1\times10^{-46}\,\mathrm{kg\,m^2}$。
代入
\[
\boxed{\omega_{\rm rms}\simeq
6.3\times10^{12}\,\mathrm{rad/s}}.
\]
\excheck 不能把两转动自由度漏成一项，
否则根号内系数将差一倍。

\thermokineticproblem{79}{72}
刚性双原子气体最初处于标准状况，
然后体积被可逆绝热压缩
$\eta=5.0$ 倍。
求末态单个分子的平均转动动能。
\exsolution
刚性双原子气体
$\gamma=7/5$。
绝热压缩至 $V_f=V_0/\eta$，
温度为
\[
T_f=T_0\left(\frac{V_0}{V_f}\right)^{\gamma-1}
=T_0\eta^{2/5}.
\]
单分子的两转动自由度
共贡献平均动能 $k_BT_f$，
所以
\[
\boxed{\overline\varepsilon_{\rm rot}
=k_BT_0\eta^{2/5}}.
\]
取 $T_0=273.15\,\mathrm K$、$\eta=5$：
\[
\boxed{\overline\varepsilon_{\rm rot}
\simeq7.2\times10^{-21}\,\mathrm J
\approx0.7\times10^{-20}\,\mathrm J}.
\]
\excheck 旋转动能为 $k_BT$，
而非平动能 $3k_BT/2$。

\thermokineticproblem{80}{72}
由刚性双原子分子组成的气体，
可逆绝热膨胀到初体积的 $\eta$ 倍。
问单位壁面积每单位时间受到分子碰撞的频率
减小多少倍？
\exsolution
分子对平面壁的碰撞通量
$\Phi=\frac14 n\bar v$，
其中数密度 $n\propto1/V$，
平均速率 $\bar v\propto\sqrt T$。
可逆绝热过程
$T\propto V^{-(\gamma-1)}$，
因此
\[
\Phi\propto
V^{-1}V^{-(\gamma-1)/2}
=V^{-(\gamma+1)/2}.
\]
刚性双原子气体
$\gamma=7/5$，
\[
\boxed{\frac{\Phi_1}{\Phi_2}
=\eta^{(\gamma+1)/2}
=\eta^{6/5}}.
\]
\excheck 既考虑了膨胀导致的分子数密度降低，
也考虑了温度下降使平均分子速率降低。

\thermokineticproblem{81}{72}
刚性双原子气体沿摩尔过程热容
$C=R$ 的多方路径，体积增大
$\eta=2.0$ 倍。
问单位面积壁面碰撞频率减少多少倍？
\exsolution
对每摩尔理想气体有
$\delta Q=C\,\mathrm dT=
C_V\,\mathrm dT+p\,\mathrm dV$。
刚性双原子气体
$C_V=5R/2$，所以
\[
p\,\mathrm dV=(C-C_V)\mathrm dT
=-\frac{3R}{2}\mathrm dT.
\]
用 $pV=RT$，
$\mathrm dT/T=-(2/3)\mathrm dV/V$，
从而 $T\propto V^{-2/3}$。
碰撞频率
$\Phi\propto n\bar v\propto V^{-1}T^{1/2}
\propto V^{-4/3}$，故
\[
\boxed{\frac{\Phi_1}{\Phi_2}
=\eta^{4/3}=2^{4/3}\simeq2.52}.
\]
\excheck 相同体积变化的不同多方路径
不一定具有与绝热过程相同的碰撞频率比。

\thermokineticproblem{82}{72}
刚性双原子气体沿某多方过程膨胀，
但分子撞击容器壁的单位面积频率
保持不变。求该过程的摩尔热容。
\exsolution
碰撞通量
$\Phi=\tfrac14n\bar v\propto V^{-1}\sqrt T$。
由于 $\Phi$ 不变，
$\sqrt T/V=\text{常数}$，
即
\[
T\propto V^2,\qquad
\frac{\mathrm dT}{T}
=2\frac{\mathrm dV}{V}.
\]
对一摩尔气体
$p=RT/V$，所以
\[
p\,\mathrm dV=\frac R2\,\mathrm dT.
\]
根据第一定律，
\[
C=C_V+\frac{p\,\mathrm dV}{\mathrm dT}
=\frac{5R}{2}+\frac R2,
\quad
\boxed{C=3R}.
\]
\excheck 频率不变并不等价于温度不变；
膨胀使数密度下降，必须提高温度以补偿碰撞速率。

\thermokineticproblem{83}{72}
在标准大气压强 $p_0$ 下，
气体质量密度为
$\rho=1.00\,\mathrm{g/L}$。
求单分子速率分布的
最概然速率、平均速率和均方根速率。
\exsolution
理想气体给
$p_0=nk_BT$、$\rho=nm$，
故
$k_BT/m=p_0/\rho$。
麦克斯韦分布的三个特征速率为
\[
\boxed{v_p=\sqrt{\frac{2p_0}{\rho}}},\qquad
\boxed{\bar v=\sqrt{\frac{8p_0}{\pi\rho}}},\qquad
\boxed{v_{\rm rms}=\sqrt{\frac{3p_0}{\rho}}}.
\]
将
$p_0=1.01325\times10^5\,\mathrm{Pa}$、
$\rho=1.00\,\mathrm{kg/m^3}$ 代入：
\[
\boxed{v_p\simeq0.45\,\mathrm{km/s}},\quad
\boxed{\bar v\simeq0.51\,\mathrm{km/s}},\quad
\boxed{v_{\rm rms}\simeq0.55\,\mathrm{km/s}}.
\]
\excheck 三者次序
$v_p<\bar v<v_{\rm rms}$，
并不需要预先知道该气体的摩尔质量。
